% =====================================================================
\section*{Перед первым уроком: сигнал, программа и безопасность}
\addcontentsline{toc}{section}{Перед первым уроком: сигнал, программа и безопасность}
\markboth{Перед первым уроком: сигнал, программа и безопасность}{}
\renewcommand{\thelisting}{В.\arabic{listing}}

Эта страница --- фундамент всей книги. Не нужно заранее знать физику или C++. Сейчас достаточно
понять одну идею: плата получает информацию, \emph{измеряя напряжение}, и действует, \emph{меняя
напряжение} на своих выводах.

\begin{tip}
Примеры этой главы запускаются после установки среды в уроке~0. Если Arduino IDE ещё не
установлена, прочитайте главу сейчас, а примеры выполните сразу после урока~0.
\end{tip}

\subsection*{Что такое электрический сигнал}

У ESP32-S3 есть выводы GPIO --- маленькие контакты, к которым подключают кнопку, светодиод или датчик.
Напряжение на выводе всегда сравнивают с \textbf{GND} (землёй, нулём схемы). Поэтому у двух
соединённых устройств обязательно должна быть общая земля.

\begin{center}
\begin{tikzpicture}[node distance=12mm]
  \node[blk,fill=blkcpu,minimum width=30mm] (chip) {ESP32-S3};
  \node[blk,fill=blkper,minimum width=30mm,left=20mm of chip] (button) {Кнопка\\ \scriptsize меняет сигнал};
  \node[blk,fill=blkper,minimum width=30mm,right=20mm of chip] (led) {Светодиод\\ \scriptsize показывает сигнал};
  \draw[arr] (button) -- node[above,font=\scriptsize,align=center]{GPIO:\\ 0 или 3,3 В} (chip);
  \draw[arr] (chip) -- node[above,font=\scriptsize,align=center]{GPIO:\\ 0 или 3,3 В} (led);
  \draw[thick,gray] (button.south) -- ++(0,-7mm) -| node[pos=.25,below,font=\scriptsize]{общий GND} (led.south);
\end{tikzpicture}
\end{center}

\textbf{Цифровой} сигнал имеет два понятных состояния: LOW --- около 0 В, HIGH --- около 3,3 В.
Позже потенциометр покажет \textbf{аналоговый} сигнал: там напряжение может быть любым между этими
границами. GPIO ESP32-S3 рассчитаны только на 3,3 В: 5 В на GPIO может необратимо повредить чип.

\subsection*{Как программа управляет платой}

Скетч Arduino --- это текст с двумя обязательными функциями. \code{void} означает: функция ничего
не возвращает. \code{setup()} запускается один раз после включения. \code{loop()} затем повторяется
снова и снова, пока плата получает питание.

\codecap{Самая маленькая программа: увидеть работу \code{setup()} и \code{loop()}}
\codeinput{l0_setup_loop}

Откройте Serial Monitor на скорости 115200. Слово «Старт» появится один раз, а строка из \code{loop()}
--- много раз. Это и есть жизненный цикл простой программы.

\subsection*{Первый сигнал без подключения деталей}

Команда ниже переводит GPIO4 в режим выхода, затем задаёт на нём HIGH и LOW. Светодиод пока не нужен:
мы проверяем логику в Serial Monitor. В уроке 1 этот же сигнал зажжёт настоящий LED через резистор.

\codecap{Вывод GPIO меняет цифровой сигнал}
\codeinput{l0_pin_signal}

\begin{caution}
Никогда не соединяйте GPIO напрямую с 5 В, с мотором, реле или светодиодной лентой. GPIO --- это
сигнальный вывод, а не источник питания. Светодиоду нужен резистор; мощной нагрузке --- драйвер или
транзисторный ключ. До любой новой сборки отключайте USB и сначала сверяйте схему.
\end{caution}

\begin{summary}
\begin{itemize}[leftmargin=*]
  \item Сигнал --- это напряжение относительно GND; цифровой сигнал ESP32-S3: LOW \(\approx\) 0 В или HIGH \(\approx\) 3,3 В.
  \item \code{setup()} работает один раз, \code{loop()} повторяется; сначала учимся наблюдать результат через Serial Monitor.
  \item Пины GPIO безопасны только для 3,3 В и малых токов.
\end{itemize}
\end{summary}

\begin{tasks}
\begin{enumerate}
\item \lv{1} Измените сообщение в \code{setup()} на своё имя, а сообщение в \code{loop()} --- на номер текущего шага.
\item \lv{1} На схеме платы найдите GND, 3V3, GPIO4, BOOT и RST (EN). Пока ничего к ним не подключайте.
\item \lv{2} Поменяйте задержку в \code{loop()} на 1000 мс. Что изменится в Serial Monitor? Сколько строк придёт за 10 секунд?
\item \lv{2} Добавьте в \code{loop()} счётчик и выводите «Шаг 1», «Шаг 2» и так далее. Почему переменная должна быть объявлена вне функции?
\item \lv{3} Выведите в Serial Monitor «таблицу» из трёх строк, повторяющуюся каждые 2 секунды: имя, номер шага и время работы платы в секундах (\code{millis() / 1000}).
\item \lv{4} Сделайте так, чтобы программа раз в секунду печатала «HIGH», а в следующую секунду «LOW», имитируя сигнал на выводе, и нарисуйте его график на бумаге.
\end{enumerate}
\end{tasks}

% дальше листинги нумеруются по урокам
\renewcommand{\thelisting}{\thesection.\arabic{listing}}
